\chapter{Proof Scope and Repository Structure}
\label{ch:proof-scope}

\section{Separated Dissertation and Proof Folders}

The repository now contains a dissertation folder with the same structure as the
reference NTT dissertation:

\begin{lstlisting}[language=bash]
dissertation/
  main.tex
  Makefile
  bibliography/refs.bib
  chapters/
  figures/README.txt
  record.md
\end{lstlisting}

The proof itself is managed separately under:

\begin{lstlisting}[language=bash]
provable-security/
  proof-files.txt
  verify-provable-security.sh
  easycrypt/
  logs/
\end{lstlisting}

The dissertation folder explains the proof.  The \code{provable-security/}
folder defines and checks the proof surface.

\section{Supplemental Material}

Material that is useful for traceability but too mechanical for the main
argument is stored under:

\begin{lstlisting}[language=bash]
dissertation/etc/
\end{lstlisting}

This includes the exhaustive generated individual EasyCrypt correspondence
appendix.  The main dissertation keeps curated proof-object analyses and
mathematical ledgers; the supplemental folder preserves the generated audit
material without forcing the reader through thousands of templated entries.

\section{Proof Manifest}

The manifest lists the EasyCrypt files that constitute the provable-security
proof gate.  It includes the main theorem layer, generic signature games, hop
games, reductions, ROM programming, scheme and transcript definitions,
distribution and algebra support, parameters, and modeled hash-interface support.

\section{EasyCrypt File Relationships}
\label{sec:easycrypt-file-relationships}

The files in the manifest form a directed proof-dependency graph.  In
Table~\ref{tab:easycrypt-file-relationships}, an entry \(F \leftarrow G\)
means that \(F\) imports \(G\) directly through an EasyCrypt
\code{require import} declaration.  Standard EasyCrypt library imports such as
\code{AllCore}, \code{Real}, \code{List}, \code{Distr}, and \code{PROM} are
omitted; Figure~\ref{fig:easycrypt-dependency-graph} and the table record only
relationships between manifest files.

Read from the leaves upward, the graph has five layers.  First,
\code{HAETAE\_Keccak1600.ec}, \code{HAETAE\_FIPS202.ec}, and
\code{HAETAE\_Params.ec} provide the hash and parameter foundations.  Second,
\code{HAETAE\_Algebra.ec}, \code{HAETAE\_Distributions.ec},
\code{HAETAE\_Assumptions.ec}, and \code{HAETAE\_Events.ec} define the
algebraic, probabilistic, hardness-relation, and event vocabulary.  Third,
\code{Sig\_ROM.ec}, \code{HAETAE\_ROM.ec}, \code{HAETAE\_Scheme.ec},
\code{HAETAE\_Transcript.ec}, \code{HAETAE\_Rejection.ec}, and
\code{HAETAE\_ROM\_Programming.ec} build the signature, transcript, rejection,
and random-oracle programming interfaces.  Fourth,
\code{HAETAE\_HopGames.ec} proves the hybrid and simulation relationships.
Finally, \code{HAETAE\_Security.ec} imports the
supporting layers and assembles the correctness and \eufcma{} theorem family.

\begin{figure}[p]
\centering
\resizebox{\textwidth}{!}{%
\begin{tikzpicture}[
  file/.style={
    rectangle,
    rounded corners=2pt,
    draw=black!55,
    line width=0.35pt,
    align=center,
    inner xsep=5pt,
    inner ysep=4pt,
    minimum width=2.65cm,
    minimum height=0.72cm,
    font=\scriptsize\ttfamily
  },
  rootfile/.style={file, fill=BrickRed!10, draw=BrickRed!70!black},
  hopfile/.style={file, fill=NavyBlue!8, draw=NavyBlue!70!black},
  interfacefile/.style={file, fill=OliveGreen!8, draw=OliveGreen!70!black},
  supportfile/.style={file, fill=Goldenrod!12, draw=Goldenrod!70!black},
  basefile/.style={file, fill=gray!10, draw=gray!70!black},
  dep/.style={
    -{Stealth[length=2.1mm,width=1.5mm]},
    draw=black!38,
    line width=0.32pt,
    shorten >=1.5pt,
    shorten <=1.5pt
  },
  spine/.style={
    -{Stealth[length=2.3mm,width=1.7mm]},
    draw=black!72,
    line width=0.55pt,
    shorten >=1.5pt,
    shorten <=1.5pt
  }
]
  \node[rootfile] (security) at (0,0) {HAETAE\_Security.ec};

  \node[hopfile] (hopgames) at (0,-1.45) {HAETAE\_HopGames.ec};

  \node[interfacefile] (sig) at (-7.3,-3.05) {Sig\_ROM.ec};
  \node[interfacefile] (scheme) at (-4.35,-3.05) {HAETAE\_Scheme.ec};
  \node[interfacefile] (romprog) at (-1.25,-3.05) {HAETAE\_ROM\\\_Programming.ec};
  \node[interfacefile] (rejection) at (2.0,-3.05) {HAETAE\_Rejection.ec};
  \node[interfacefile] (transcript) at (5.2,-3.05) {HAETAE\_Transcript.ec};
  \node[interfacefile] (rom) at (8.25,-3.05) {HAETAE\_ROM.ec};

  \node[supportfile] (reductions) at (-5.9,-4.9) {HAETAE\_Reductions.ec};
  \node[supportfile] (events) at (-2.7,-4.9) {HAETAE\_Events.ec};
  \node[supportfile] (assumptions) at (0.75,-4.9) {HAETAE\_Assumptions.ec};
  \node[supportfile] (distributions) at (4.25,-4.9) {HAETAE\\\_Distributions.ec};

  \node[basefile] (algebra) at (-2.65,-6.55) {HAETAE\_Algebra.ec};
  \node[basefile] (params) at (1.15,-6.55) {HAETAE\_Params.ec};
  \node[basefile] (fips) at (5.05,-6.55) {HAETAE\_FIPS202.ec};
  \node[basefile] (keccak) at (5.05,-8.05) {HAETAE\\\_Keccak1600.ec};

  \draw[spine] (security) -- (hopgames);
  \draw[spine] (security) -- (romprog);
  \draw[spine] (security) -- (rejection);
  \draw[spine] (security) -- (transcript);
  \draw[spine] (security) -- (reductions);

  \draw[dep] (security) -- (sig);
  \draw[dep] (security) -- (scheme);
  \draw[dep] (security) -- (algebra);
  \draw[dep] (security) -- (distributions);
  \draw[dep] (security) -- (assumptions);

  \draw[dep] (hopgames) -- (sig);
  \draw[dep] (hopgames) -- (scheme);
  \draw[dep] (hopgames) -- (params);
  \draw[dep] (hopgames) -- (algebra);
  \draw[dep] (hopgames) -- (distributions);
  \draw[dep] (hopgames) -- (assumptions);
  \draw[dep] (hopgames) -- (transcript);
  \draw[dep] (hopgames) -- (reductions);
  \draw[dep] (hopgames) -- (rejection);
  \draw[dep] (hopgames) -- (rom);
  \draw[dep] (hopgames) -- (romprog);

  \draw[dep] (romprog) -- (params);
  \draw[dep] (romprog) -- (algebra);
  \draw[dep] (romprog) -- (distributions);
  \draw[dep] (romprog) -- (rom);
  \draw[dep] (romprog) -- (transcript);
  \draw[dep] (romprog) -- (rejection);
  \draw[dep] (romprog) -- (reductions);

  \draw[dep] (scheme) -- (sig);
  \draw[dep] (scheme) -- (params);
  \draw[dep] (scheme) -- (algebra);
  \draw[dep] (scheme) -- (distributions);
  \draw[dep] (scheme) -- (rom);

  \draw[dep] (transcript) -- (params);
  \draw[dep] (transcript) -- (algebra);
  \draw[dep] (transcript) -- (rom);
  \draw[dep] (transcript) -- (assumptions);

  \draw[dep] (rejection) -- (params);
  \draw[dep] (rejection) -- (algebra);
  \draw[dep] (rejection) -- (distributions);
  \draw[dep] (rejection) -- (events);
  \draw[dep] (rejection) -- (reductions);
  \draw[dep] (rejection) -- (transcript);

  \draw[dep] (rom) -- (params);
  \draw[dep] (rom) -- (algebra);
  \draw[dep] (rom) -- (distributions);

  \draw[dep] (reductions) -- (events);
  \draw[dep] (events) -- (params);
  \draw[dep] (events) -- (algebra);
  \draw[dep] (assumptions) -- (params);
  \draw[dep] (assumptions) -- (algebra);
  \draw[dep] (distributions) -- (params);
  \draw[dep] (distributions) -- (algebra);
  \draw[dep] (algebra) -- (params);
  \draw[dep] (algebra) -- (fips);
  \draw[dep] (fips) -- (keccak);
\end{tikzpicture}%
}
\caption{Manifest-level dependency graph for the HAETAE EasyCrypt proof.  An
arrow points from an importing file to a directly imported manifest file; the
darker arrows highlight the main top-level assembly paths.}
\label{fig:easycrypt-dependency-graph}
\end{figure}

\begin{longtable}{p{0.30\textwidth}p{0.58\textwidth}}
\caption{Immediate manifest-file relationships induced by EasyCrypt imports.}
\label{tab:easycrypt-file-relationships}\\
\toprule
\textbf{File} & \textbf{Direct manifest imports} \\
\midrule
\endfirsthead
\toprule
\textbf{File} & \textbf{Direct manifest imports} \\
\midrule
\endhead
\code{HAETAE\_Security.ec} &
\code{Sig\_ROM.ec}, \code{HAETAE\_Scheme.ec}, \code{HAETAE\_Algebra.ec},
\code{HAETAE\_Distributions.ec}, \code{HAETAE\_Reductions.ec},
\code{HAETAE\_Assumptions.ec}, \code{HAETAE\_Rejection.ec},
\code{HAETAE\_Transcript.ec},
\code{HAETAE\_ROM\_Programming.ec}, \code{HAETAE\_HopGames.ec}. \\
\midrule
\code{Sig\_ROM.ec} &
None inside the manifest; it depends only on EasyCrypt library files. \\
\midrule
\code{HAETAE\_HopGames.ec} &
\code{Sig\_ROM.ec}, \code{HAETAE\_Scheme.ec}, \code{HAETAE\_Params.ec},
\code{HAETAE\_Algebra.ec}, \code{HAETAE\_Distributions.ec},
\code{HAETAE\_Assumptions.ec}, \code{HAETAE\_Transcript.ec},
\code{HAETAE\_Reductions.ec}, \code{HAETAE\_Rejection.ec},
\code{HAETAE\_ROM.ec}, \code{HAETAE\_ROM\_Programming.ec}. \\
\midrule
\code{HAETAE\_Reductions.ec} &
\code{HAETAE\_Events.ec}. \\
\midrule
\code{HAETAE\_ROM.ec} &
\code{HAETAE\_Params.ec}, \code{HAETAE\_Algebra.ec},
\code{HAETAE\_Distributions.ec}. \\
\midrule
\code{HAETAE\_ROM\_Programming.ec} &
\code{HAETAE\_Params.ec}, \code{HAETAE\_Algebra.ec},
\code{HAETAE\_Distributions.ec}, \code{HAETAE\_ROM.ec},
\code{HAETAE\_Transcript.ec}, \code{HAETAE\_Rejection.ec},
\code{HAETAE\_Reductions.ec}. \\
\midrule
\code{HAETAE\_Scheme.ec} &
\code{Sig\_ROM.ec}, \code{HAETAE\_Params.ec}, \code{HAETAE\_Algebra.ec},
\code{HAETAE\_Distributions.ec}, \code{HAETAE\_ROM.ec}. \\
\midrule
\code{HAETAE\_Transcript.ec} &
\code{HAETAE\_Params.ec}, \code{HAETAE\_Algebra.ec},
\code{HAETAE\_ROM.ec}, \code{HAETAE\_Assumptions.ec}. \\
\midrule
\code{HAETAE\_Events.ec} &
\code{HAETAE\_Params.ec}, \code{HAETAE\_Algebra.ec}. \\
\midrule
\code{HAETAE\_Assumptions.ec} &
\code{HAETAE\_Params.ec}, \code{HAETAE\_Algebra.ec}. \\
\midrule
\code{HAETAE\_Distributions.ec} &
\code{HAETAE\_Params.ec}, \code{HAETAE\_Algebra.ec}. \\
\midrule
\code{HAETAE\_Rejection.ec} &
\code{HAETAE\_Params.ec}, \code{HAETAE\_Algebra.ec},
\code{HAETAE\_Distributions.ec}, \code{HAETAE\_Events.ec},
\code{HAETAE\_Reductions.ec}, \code{HAETAE\_Transcript.ec}. \\
\midrule
\code{HAETAE\_Algebra.ec} &
\code{HAETAE\_Params.ec}, \code{HAETAE\_FIPS202.ec}. \\
\midrule
\code{HAETAE\_Params.ec} &
None inside the manifest; it depends only on EasyCrypt library files. \\
\midrule
\code{HAETAE\_FIPS202.ec} &
\code{HAETAE\_Keccak1600.ec}. \\
\midrule
\code{HAETAE\_Keccak1600.ec} &
None inside the manifest; it depends only on EasyCrypt library files. \\
\bottomrule
\end{longtable}

\section{Theorem Dependency Classes}

To avoid overstating the proof, the dissertation classifies every final security
claim into one of four dependency classes.

\begin{longtable}{p{0.20\textwidth}p{0.68\textwidth}}
\caption{Dependency classes used for theorem-scope statements.}
\label{tab:dependency-classes}\\
\toprule
\textbf{Class} & \textbf{Meaning} \\
\midrule
\endfirsthead
\toprule
\textbf{Class} & \textbf{Meaning} \\
\midrule
\endhead
Checked & A statement accepted by EasyCrypt for the files in
\code{provable-security/proof-files.txt}. \\
\midrule
Assumed & A named modeling hypothesis needed to read the checked EasyCrypt model
as the intended HAETAE specification model. \\
\midrule
Audited & A specification-to-model correspondence supported by local inspection
and tables, but not by a checked equivalence theorem. \\
\midrule
Out of scope & A claim deliberately not proved by the EasyCrypt security
development, such as concrete SHAKE standard-model security, implementation
parsing correctness, concrete MLWE/Module-SIS bit estimates, or reduction
runtime bounds. \\
\bottomrule
\end{longtable}

The central theorem is in the checked class only for the EasyCrypt formal model.
The specification-level reading is conditional on the assumed and audited
classes.  This dependency discipline is part of the proof artifact: it prevents
a verified game proof from being silently promoted into a stronger end-to-end
certification claim.

\section{Excluded Generated KAT Artifact}

\code{HAETAE\_FIPS202\_KAT\_Empty\_Generated.ec} is not part of the
provable-security gate.  It is a generated FIPS202 known-answer-test artifact.
It may be useful for implementation validation or certification evidence, but
it is not a reduction step in the \eufcma{} proof.  This distinction prevents a
large generated file from obscuring the status of the security proof.

The same distinction applies to modeled hash-interface support more generally.  The
security theorem is a random-oracle-model theorem for the formal model.  FIPS202
and Keccak support files help define or validate modeled hash objects, but they
do not by themselves prove standard-model security of a concrete SHAKE
instantiation.

\section{ROM-to-FIPS202 Boundary}

The proof uses typed random-oracle queries and lazy-ROM laws to reason about
freshness, programming, and adversarial prequeries.  The FIPS202 and Keccak
files are therefore not a replacement for the ROM theorem; they occupy a
different verification layer.  They may support implementation conformance or
hash-interface documentation, but the security theorem continues to assume the
random-oracle abstraction.  Any theorem replacing the ROM by SHAKE would need an
additional assumption or proof principle, for example an indifferentiability or
computational random-oracle-instantiation argument.  No such theorem is part of
the checked provable-security gate described here.
